% --- SECCIÓN 4: DISEÑO Y PROTOTIPADO ---
\section{Diseño y Prototipado}
\label{sec:diseno_prototipado}

La traducción visual de EthosHub busca eliminar la apariencia de "prototipo estudiantil", aplicando estándares de diseño de interfaces corporativas.

\subsection{Arquitectura de Información}
Se dividió la navegación en dos áreas excluyentes:
\begin{itemize}[noitemsep]
    \item \textbf{Zona Privada (Gestión):} Formularios y dashboards accesibles solo mediante validación de token JWT.
    \item \textbf{Zona Pública (Proyección):} Renderizado de solo lectura del portafolio profesional.
\end{itemize}

\subsection{Mockups de Alta Fidelidad}
Los prototipos aprobados para el inicio del desarrollo incluyen:

\begin{itemize}
    \item \textbf{Panel de Autenticación:} Formularios limpios para Login/Registro, preparados para integrar flujos OAuth2.
    \item \textbf{Dashboard (Epic 1):} Vista central para la edición del perfil, utilizando \comando{JetBrains Mono} en las áreas de configuración técnica o despliegue de repositorios.
    \item \textbf{Portafolio Público (Epic 2):} Diseño de \textit{Single Page} optimizado para lectura rápida por parte de reclutadores IT.
\end{itemize}

\begin{table}[!ht]
    \centering
    \begin{tabularx}{\textwidth}{|l|X|l|}
        \hline
        \rowcolor{colorPrimario} \textcolor{white}{\textbf{Módulo Visual}} & \textcolor{white}{\textbf{Foco de UX}} & \textcolor{white}{\textbf{Estado}} \\ \hline
        Auth (Login/Register) & Fricción mínima al ingreso. & \estadoPass \\ \hline
        Profile Manager & Componentización mediante \textit{Hooks}. & \estadoPass \\ \hline
        Public View & Jerarquía de tipografía \comando{Sora}. & \estadoPass \\ \hline
    \end{tabularx}
    \caption{Estado de los Prototipos UI.}
\end{table}

\begin{NotaTecnica}
    Los Mockups fueron creados en Figma y se encuentran disponibles en el siguiente enlace: \url{https://url-shortener.me/IT5L}
\end{NotaTecnica}